\begin{proof}[Proof of \cref{thm:p143-dominated-divisorial}]
We use analytic resolution and principalization of coherent ideals,
together with the ordinary divisorial-density theorem
\cite[Theorem~2.3.6]{MP24}. The latter will only be used to lift a
semivaluation through a modification. The order inequality will be
proved separately.

\medskip\noindent
\emph{Step 1: a lift exists, even when the support is proper.}
Let \(\pi:Y\to X\) be a proper bimeromorphic morphism of connected
compact K\"ahler manifolds. Pullback of coherent ideals induces a
continuous map \(\pi^\beth:Y^\beth\to X^\beth\).
The space \(Y^\beth\) is compact and \(X^\beth\) is Hausdorff, so its
image is closed. Every point of \(X^{\mathrm{div}}\) is in this image
by birational invariance of divisorial valuations
\cite[Remark~2.3.2]{MP24}. Ordinary divisorial density therefore gives
\[
                           \pi^\beth(Y^\beth)=X^\beth.
\]
We can consequently choose \(\widetilde v\in Y^\beth\) lifting any
given \(v\in X^\beth\). We assert existence, not uniqueness. In
particular we have not extended the valuation-only bijection in
\cite[Proposition~2.1.6]{MP24} by an undefined
\(+\infty-(+\infty)\) formula.

\medskip\noindent
\emph{Step 2: a rational monomial valuation dominated by a semivaluation.}
Suppose that \(B_1,\ldots,B_k\) are distinct smooth prime divisors
with simple normal crossings on \(Y\), that
\(\widetilde v(B_i)>0\), and that
\[
                0<a_i\leq\widetilde v(B_i),\qquad a_i\in\mathbf Q .
\]
Infinite values of \(\widetilde v(B_i)\) are allowed.
Here \(\widetilde v(B_i)\) means the value of
\(\mathcal O_Y(-B_i)\).

Let \(Z=Z_Y(\widetilde v)\) be the irreducible central variety supplied
by \cite[Lemma~2.1.1]{MP24}. The ideal defining \(Z\) is maximal among
coherent ideals with positive value. Thus \(Z\subset B_i\) for every
\(i\). There is a unique connected component \(W\) of
\(B_1\cap\cdots\cap B_k\) containing \(Z\): the intersection is smooth
and its connected components are disjoint. Let \(q\) be the monomial
valuation along \(W\) with weights \((a_1,\ldots,a_k)\).
The construction and independence of the point of \(W\) are
\cite[Section~2.3.1, Proposition~2.3.4]{MP24}.
Rationality of the weights makes \(q\) divisorial, with the rational
scaling allowed here; see \cite[Remark~2.3.5]{MP24}.

We claim, importantly on \emph{every} coherent ideal of \(Y\), that
\begin{equation}
                        q(\mathfrak j)\leq\widetilde v(\mathfrak j).
 \label{eq:p143-monomial-domination}
\end{equation}
For \(s>0\), form the coherent monomial ideal
\begin{equation}
 \mathfrak c_s=
 \sum_{\substack{\beta\in\mathbf Z_{\geq0}^k\\
                  \sum_i a_i\beta_i\geq s}}
       \mathcal O_Y\left(-\sum_i\beta_iB_i\right).
 \label{eq:p143-monomial-ideal}
\end{equation}
This is a finite sum after discarding redundant generators.
Indeed the exponent set is upward closed; a minimal exponent cannot
have \(\beta_i>\lceil s/a_i\rceil\). The semivaluation rules therefore
give, in extended-real order,
\begin{equation}
 \widetilde v(\mathfrak c_s)
  =\min_\beta\sum_i\beta_i\widetilde v(B_i)
  \geq\min_\beta\sum_i\beta_i a_i\geq s .
 \label{eq:p143-monomial-ideal-bound}
\end{equation}
Terms with exponent zero do not contribute to these sums, so no
undefined product \(0\cdot(+\infty)\) occurs.

Take a nonzero coherent ideal \(\mathfrak j\), and suppose
\(0<s<q(\mathfrak j)\). In coordinates near any point of \(W\), with
\(B_i=(z_i=0)\), the ideal \(\mathfrak c_s\) is generated by
the monomials \(z^\beta\) of weight at least \(s\).
A holomorphic germ has monomial value at least \(s\) precisely when
all its transverse Taylor coefficients of weight less than \(s\)
vanish identically as germs along \(W\). There are only finitely many
such transverse exponents, so this is equivalent to membership in the
displayed monomial ideal. The coordinates tangent to \(W\) have weight
zero; vanishing at one point of \(W\) is not the same as vanishing of
one of these coefficient germs. By the point-independence of the
monomial valuation and local finite generation of \(\mathfrak j\),
we obtain
\[
                  \mathfrak j\subseteq\mathfrak c_s
                  \quad\hbox{on a neighbourhood of }W.
\]
The inclusion need not hold globally, because the intersection of the
\(B_i\) can have other connected components. To make the argument
global, set
\[
 \mathcal Q=(\mathfrak j+\mathfrak c_s)/\mathfrak c_s,\qquad
 \mathfrak k=\operatorname{Ann}_{\mathcal O_Y}(\mathcal Q).
\]
Both are coherent. Since \(\mathcal Q\) vanishes near \(W\),
\(\mathfrak k=\mathcal O_Y\) there, and
\(\mathfrak k\mathfrak j\subseteq\mathfrak c_s\) globally.
We have \(\widetilde v(\mathfrak k)=0\): a positive value would imply
\(\mathfrak k\subseteq\mathcal I_Z\), contradicting that
\(\mathfrak k\) is the unit ideal near \(Z\subset W\).
Consequently
\[
 \widetilde v(\mathfrak j)
 =\widetilde v(\mathfrak k\mathfrak j)
 \geq\widetilde v(\mathfrak c_s)\geq s .
\]
Let \(s\) tend to \(q(\mathfrak j)\). If \(q(\mathfrak j)=0\) the
inequality is automatic; the zero ideal has infinite value on both
sides. This proves \eqref{eq:p143-monomial-domination}.
It also explains why generic-stratum information suffices at the
possibly smaller actual centre \(Z\).

\medskip\noindent
\emph{Step 3: approximate finitely many ideal values.}
Fix finitely many nonzero coherent ideals
\(\mathfrak a_1,\ldots,\mathfrak a_\ell\) on \(X\).
Simultaneous analytic principalization gives a proper modification
\(\pi:Y\to X\), with \(Y\) smooth, such that
\[
 \mathfrak a_j\mathcal O_Y=\mathcal O_Y(-D_j),\qquad
 D_j=\sum_{i=1}^r d_{ji}B_i,\qquad d_{ji}\in\mathbf Z_{\geq0},
\]
and the distinct \(B_i\) form a simple normal crossing divisor.
The modification may be taken as a finite sequence of blowups of
smooth centres; thus \(Y\) is again compact K\"ahler.
One can, for example, principalize their product and resolve the
resulting pair. The individual invertible factors follow locally
because a product of nonzero ideals that is principal and invertible
has invertible factors. This is the usual analytic log-resolution
input, not an assertion about reduced semistable models over a
valued field.

Choose a lift \(\widetilde v\) by Step 1, and put
\(w_i=\widetilde v(B_i)\in[0,+\infty]\). Multiplicativity gives
\begin{equation}
                         v(\mathfrak a_j)=\sum_i d_{ji}w_i.
 \label{eq:p143-principalized-values}
\end{equation}
Retain only indices with \(w_i>0\). For those indices choose finite
positive rational numbers \(a_i\leq w_i\). For the others put \(a_i=0\).
If no index is retained, use the trivial valuation. Otherwise Step 2
constructs a rational divisorial valuation \(q_Y\leq\widetilde v\),
and its pushforward \(q_X=\pi^\beth(q_Y)\) is a member of
\(X^{\mathrm{div}}\) satisfying
\[
            q_X\leq v,\qquad
            q_X(\mathfrak a_j)=\sum_i d_{ji}a_i .
\]
A component \(B_i\) whose weight is zero does not contain the generic
point of the selected stratum: otherwise simple normal crossings
would fail there. Its monomial value is therefore zero, as used in
this formula.

For each finite positive \(w_i\), let \(a_i\) approach \(w_i\) from
below through positive rationals; for \(w_i=+\infty\), let \(a_i\)
increase without bound. The finite sums then converge to
\eqref{eq:p143-principalized-values} in \([0,+\infty]\).
More concretely, for any integer \(N>0\) we can arrange
\[
 \begin{cases}
 0\leq v(\mathfrak a_j)-q_X(\mathfrak a_j)<1/N,
       &v(\mathfrak a_j)<+\infty,\\
 q_X(\mathfrak a_j)>N,&v(\mathfrak a_j)=+\infty .
 \end{cases}
\]
The zero and unit ideals already have their prescribed values.
Basic neighbourhoods of \(v\) involve only finitely many ideals.
We have therefore constructed a dominated rational divisorial point
in every such neighbourhood. This proves
\eqref{eq:p143-dominated-closure}.
For an explicit net, direct pairs consisting of a finite set of
nonzero coherent ideals and a positive integer \(N\) by inclusion
and the usual order, and make the above choice for each pair.
No order relation between two different chosen points is needed.
\end{proof}

\begin{proof}[Proof of \cref{cor:p143-usc-antitone-comparison}]
MP potentials are upper semicontinuous, being decreasing limits of
continuous PL functions. Fix \(v\), and take \(q_\lambda\leq v\)
with \(q_\lambda\to v\) from \cref{thm:p143-dominated-divisorial}.
Antitonicity and divisorial agreement imply
\[
                         F(v)\leq F(q_\lambda)=\phi(q_\lambda).
\]
Upper semicontinuity of \(\phi\), valid for nets as well as sequences,
then gives \(F(v)\leq\phi(v)\). If \(\phi(v)=-\infty\), this means
that every real upper threshold eventually bounds \(\phi(q_\lambda)\);
the same inequality still follows in extended-real order.

For the reverse inequality, apply
\cite[Theorem~4.3.8]{MP24} to the MP potential \(\phi\) and the
upper-semicontinuous function \(F\). Their equality on
\(X^{\mathrm{div}}\) gives \(\phi\leq F\) everywhere.
Combining both inequalities proves
\eqref{eq:p143-all-point-comparison}.
\end{proof}

\begin{proof}[Proof of \cref{cor:p143-finite-energy-realization}]
The finite-energy correspondence supplies
\(\phi\in\mathcal E^1_{\mathrm{MP}}([\omega])\) with
\(\phi(q)=F_\Gamma(q)\) for every rational divisorial \(q\):
use \cite[Theorem~3.17]{DXZ23} and
\cite[Theorems~5.1.4, 5.1.7 and Section~5.2]{MP24}, with the
normalization and divisorial-identification interface recorded in the
finite-energy comparison statement surrounding equations
\textup{(71.39)}--\textup{(71.40)} of the frozen Chapter~71 input.
The trivial valuation is included: both values are
\(c=\Gamma_{\max}\).

The raw transform is bounded above by \(c\). If \(q\leq v\), then
\[
 q(\mathcal I(m\Gamma_\tau))
       \leq v(\mathcal I(m\Gamma_\tau))
 \quad\hbox{for all }m,\tau.
\]
Taking the first supremum in \eqref{eq:p143-raw-transform}, changing
the sign, and taking the second supremum shows
\(F_\Gamma(q)\geq F_\Gamma(v)\).
Its upper semicontinuity is precisely the stipulated P14.2 input.
Thus \cref{cor:p143-usc-antitone-comparison} applies.

The equality concerns the literal ideal formula, including
semivaluations supported on a proper subvariety. It does not replace
that formula by an unspecified regularization. The only P14.2 input
used is its upper-semicontinuity conclusion, not positivity of its
continuous approximants.
\end{proof}
